\RequirePackage{plautopatch}
\documentclass[platex,dvipdfmx,a4paper]{jsarticle}

\usepackage{amsmath}
\usepackage{amssymb}
\newcommand{\hb}[2]{\frac{\sectionial #1}{\sectionial #2}}
\newcommand{\vectxy}[2]{\begin{pmatrix}#1\\#2\end{pmatrix}}
\newcommand{\vectxyz}[3]{\begin{pmatrix}#1\\#2\\#3\end{pmatrix}}
\newcommand{\vect}[1]{\boldsymbol{#1}}

\newcommand{\ssum}[3]{\sum_{#1=#2}^#3}

\title{物理数学I演習レポート1}
\author{嶋田 悠音}
\date{2026年5月6日}

\begin{document}
  \maketitle
  \subsection*{問題 7.1}
  
  \subsubsection*{(i) $\sin z$ の Maclaurin 展開と収束半径}
  $\sin z$ の Maclaurin 展開は、以下の級数で与えられる。
  \begin{align*}
    \sin z = \sum_{n=0}^{\infty} \frac{(-1)^n}{(2n+1)!} z^{2n+1} = z - \frac{z^3}{3!} + \frac{z^5}{5!} - \dots
  \end{align*}
  この級数の第 $n$ 項を $a_n$ とすると、ダランベールの収束判定法により、
  \begin{align*}
    \lim_{n \to \infty} \left| \frac{a_{n+1}}{a_n} \right| = \lim_{n \to \infty} \left| \frac{z^2}{(2n+2)(2n+3)} \right| = 0
  \end{align*}
  となる。これが任意の複素数 $z$ に対して成り立つため、この級数は複素平面全体で収束する。したがって、収束半径 $R$ は無限大（$R = \infty$）である。

  \subsubsection*{(ii) $1/\sin z$ の Laurent 展開の導出}
  (i) で求めた展開式を用いると、$1/\sin z$ は次のように書ける(oはランダウの記号)。
  \begin{align*}
    \frac{1}{\sin z} &= \frac{1}{z - \frac{z^3}{6} + \frac{z^5}{120} - o(z^5)} \\
    &= \frac{1}{z} \cdot \frac{1}{1 - \left( \frac{z^2}{6} - \frac{z^4}{120} + o(z^4) \right)}
  \end{align*}
  ここで、$\alpha = \frac{z^2}{6} - \frac{z^4}{120} + o(z^4)$ とおく。$|z|$ が十分に小さいとき $|\alpha| < 1$ とみなせるため、べき級数の展開公式
  \begin{align*}
    \frac{1}{1-\alpha} = 1 + \alpha + \alpha^2 + \alpha^3 + \dots
  \end{align*}
  を適用することができる。これを $z^4$ の次数まで考慮して展開すると、以下のようになる。
  \begin{align*}
    \frac{1}{\sin z} &= \frac{1}{z} \left[ 1 + \left( \frac{z^2}{6} - \frac{z^4}{120} \right) + \left( \frac{z^2}{6} - \frac{z^4}{120} \right)^2 + o(z^4) \right] \\
    &= \frac{1}{z} \left[ 1 + \frac{z^2}{6} - \frac{z^4}{120} + \left( \frac{z^4}{36} + o(z^6) \right) + 0(z^4) \right] \\
    &= \frac{1}{z} \left[ 1 + \frac{z^2}{6} + \left( \frac{-3 + 10}{360} \right) z^4 + o(z^4) \right] \\
    &= \frac{1}{z} + \frac{z}{6} + \frac{7}{360} z^3 + o(z^3)
  \end{align*}
  以上より、求める Laurent 展開の $z^3$ の項までの結果は、
  \begin{align*}
    \frac{1}{\sin z} = \frac{1}{z} + \frac{z}{6} + \frac{7}{360} z^3 + \dots
  \end{align*}
  である。

  \subsubsection*{(iii) 展開の有効範囲}
  関数 $f(z) = \frac{1}{\sin z}$ は、分母が零となる点、すなわち $\sin z = 0$ を満たす $z = n\pi$ （$n$ は整数）において孤立特異点を持つ。
  Laurent 展開 $0 < |z - z_0| < R$ が有効となる半径 $R$ は、中心 $z_0 = 0$ から最も近い特異点までの距離に等しい。
  原点以外で最も原点に近い特異点は $z = \pi$ および $z = -\pi$ である。
  したがって、この Laurent 展開が有効な範囲は、中心の極 $z=0$ を除いた以下の領域となる。
  \begin{align*}
    0 < |z| < \pi
  \end{align*}
\end{document}
